\documentclass[12pt]{article}
\usepackage[a4paper,margin=2.2cm]{geometry}
\usepackage{CJKutf8}
\usepackage{setspace}
\usepackage{parskip}

\begin{document}
\begin{CJK}{UTF8}{bsmi}
\onehalfspacing

\begin{center}
{\LARGE Recommendation Letter}
\end{center}

\vspace{1em}

\noindent 本人謹此誠摯推薦晚南申請「NCTS-Phys Research Assistant」。晚南是一位具高度責任感、研究紀律與執行力的研究生，無論在研究規劃、實際執行或成果整理上，皆能主動承擔責任並確實完成交付之研究任務，展現出成熟且可靠的研究態度。

\noindent 在研究歷程中，晚南表現出極為細心與嚴謹的特質，能敏銳觀察研究中容易被忽略的關鍵細節，並即時檢視與修正潛在問題。此一能力使他在理論分析、程式實作與數據驗證等層面，皆能維持良好的研究品質與一致性。此外，他具備優秀的學習能力，面對新理論、新工具或跨領域知識時，能迅速理解其核心概念，並進一步將其有效應用於實際研究問題之中，顯示出高度的學術適應性與成長潛力。

\noindent 在研究主題方面，晚南目前專注於近程量子電腦環境下之量子誤差緩解研究。隨著量子硬體逐步發展，實際量子裝置仍普遍受到量子閘不完美、量測誤差以及量子位元間交叉干擾等噪聲影響，使得計算結果難以維持穩定與準確。在此背景下，學生以機率式誤差抵消（Probabilistic Error Cancellation, PEC）作為研究核心，顯示其對當前量子計算發展瓶頸具有清楚且前瞻的理解。

\noindent 目前，晚南已著手開發一套適用於近程量子電腦的自動化 PEC 軟體套件，整合噪聲模型建立、量子電路生成、誤差緩解操作與結果重建流程，使研究人員能以相對簡化的設定進行小規模 PEC 模擬與測試。該套件同時支援不同噪聲模型，並朝向貼近實際量子硬體情境的應用發展，具備良好的研究延展性與實用價值，充分展現其將理論方法落實為可操作研究工具的能力。

\noindent 整體而言，晚南具備扎實的量子資訊與量子計算基礎，並展現出獨立研究能力、解決問題的耐心，以及對研究主題的高度投入與熱忱。我深信其具備良好的潛力，能勝任研究助理所需之研究執行、技術整合與團隊協作工作，並在相關研究計畫中持續學習、穩定貢獻。

\noindent 基於以上理由，本人特此誠摯推薦晚南同學申請「NCTS-Phys Research Assistant」，並對其未來在量子計算及相關跨領域研究之發展深具信心。

\vspace{2cm}

\noindent 推薦人：\underline{\hspace{4cm}}\\[0.6cm]
\noindent 職稱：\underline{\hspace{4cm}}\\[0.6cm]
\noindent 單位：\underline{\hspace{6cm}}\\[0.6cm]
\noindent 日期：\underline{\hspace{4cm}}

\end{CJK}
\end{document}
